\section{Preliminaries}\label{prelim}
We start with some useful properties of the quasi-invariants.
\begin{Proposition}[\cite{etingof2002lectures}]\label{firstProp}
 Let $\k$ be a field.
\begin{enumerate}\itemsep=0pt
 \item[$1$.] $\k[x_1,x_2,x_3]^{S_3}\subset Q_m(3,\k)$, $Q_0(3,\k) = \k[x_1,x_2,x_3]$, and $Q_m(3,\k)\supset Q_{m'}(3,\k)$, where ${m'>m}$.
 \item[$2$.] $Q_m(3,\k)$ is a ring.
 \item[$3$.] $Q_m(3,\k)$ is a finitely generated $\k[x_1,x_2,x_3]^{S_3}$-module.
\end{enumerate}
\end{Proposition}

Note that \cite{etingof2002lectures} proves Proposition~\ref{firstProp} in the case, where $\k=\C$. However, the proofs for the first two assertions work over any field, and the last assertion follows from the Hilbert basis theorem. In view of the structure of $Q_m(3,\k)$ as a module over the symmetric polynomials, given some $K\in Q_m(3,\k)$, we will frequently refer to quasi-invariant polynomials that can be obtained via scalar multiplication of $Q$ by a symmetric polynomial. To distinguish these polynomials from the ordinary $\k$-multiples of $Q$, we will refer to them as \emph{symmetric polynomial multiples} of $Q$.

We consider $Q_m(3, \mathbf{F}_{2})$ and $Q_m(3, \mathbf{F}_{3})$ as representations of $S_3$, where $S_3$ permutes the variables $x_1$, $x_2$, $x_3$. Since $Q_m(3, \mathbf{F}_{2})$ and $Q_m(3, \mathbf{F}_{3})$ are vector spaces over $\mathbf{F}_2$ and $\mathbf{F}_3$ respectively and the characteristics $2$ and $3$ divide $|S_3|$, $Q_m(3, \mathbf{F}_{2})$ and $Q_m(3, \mathbf{F}_{3})$ are modular representations~of~$S_3$.

\begin{Proposition}
 $Q_m(3, \mathbf{F}_{2})$ and $Q_m(3, \mathbf{F}_{3})$ are modular representations of $S_3$.
\end{Proposition}

First, we consider characteristic $2$.

\subsection[Preliminary definitions for p=2]{Preliminary definitions for $\boldsymbol{p=2}$}

We describe the indecomposable and irreducible representations of $S_3$ for $p=2$.

\begin{Proposition}[\cite{Alperin_1986}]\label{prop:p2indecomp}
 There are $3$ irreducible or indecomposable representations of $S_{3}$ in characteristic $2$:
 \begin{enumerate}\itemsep=0pt
 \item[$1$.] $\triv$ is the irreducible representation of $S_3$ that is acted on trivially by $S_{3}$.
 \item[$2$.] $\std$ is the $2$-dimensional irreducible representation of $S_3$ obtained by reducing the standard representation in characteristic $0$ mod~$2$.
 \item[$3$.] $\trtr$ is the $2$-dimensional indecomposable representation that contains a copy of $\triv$ as a subrepresentation such that the quotient of $\trtr$ by this subrepresentation is $\triv$.
 \end{enumerate}
\end{Proposition}

\begin{Example}\label{eg:p2reps}
 The polynomial $E_{\trtr}:=x_1^2x_2+x_2^2x_3+x_3^2x_1\in\mathbf{F}_2[x_1,x_2,x_3]$ generates a~copy of $\trtr$. To see this, note that for any $i_1$, $i_2$, we have
 \[(1-s_{i_1i_2})E_\trtr=x_1^2x_2+x_1x_2^2+x_1^2x_3+x_1x_3^2+x_2^2x_3+x_2x_3^2\in\mathbf{F}_2[x_1,x_2,x_3]^{S_3}.\]
 Since the transpositions generate $S_3$, $E_\trtr$ generates a two-dimensional representation that contains $\triv$ as a subrepresentation. Moreover, since $E_\trtr$ is not symmetric, this representation is not $\triv\oplus\triv$, so it must be $\trtr$.
\end{Example}

We then study the behaviors of each indecomposable representation in the quasi-invariants. We define $Q_m(3,\mathbf{F}_2)_{\triv}$ and $Q_m(3,\mathbf{F}_2)_{\std}$ to be the direct sum of all copies of $\triv$ and $\std$ respectively in the quasi-invariants. We also define $Q_m(3,\mathbf{F}_2)_\trtr$ to be the direct sum of all copies of $\triv$ and $\trtr$.

\begin{Remark}
 We cannot define $Q_m(3,\mathbf{F}_2)_\trtr$ to exclude copies of $\triv$ since we can add elements of $Q_m(3,\mathbf{F}_2)_\triv$ to copies of $\trtr$ and still obtain a copy of $\trtr$. For example, \smash{$F:=E_\trtr+x_1^3+x_2^3+x_3^3$} still satisfies $(1-s_{i_1i_2})F=(1-s_{i_1i_2})E_\trtr$ for all~$i_1$,~$i_2$, so it generates a copy of $\trtr$ by the same argument as Example~\ref{eg:p2reps}.
\end{Remark}

\begin{Proposition}[{\cite{wang2022explicit}}]\label{prop:p2triv}
 As an $\mathbf{F}_2[x_1,x_2,x_3]^{S_3}$-module, $Q_m(3,\mathbf{F}_2)_{\triv}$ is freely generated by $1$.
\end{Proposition}

Note that by the classification of indecomposables in Proposition~\ref{prop:p2indecomp}, every extension of $\std$ and every extension of a module by $\std$ splits. Thus $Q_m(3,\mathbf{F}_2)_\std$ is a direct summand of $Q_m(3,\mathbf{F}_2)$ (whose complement is $Q_m(3,\mathbf{F}_2)_\trtr$), and we mainly consider $Q_m(3,\mathbf{F}_2)_\std$. $Q_m(3,\mathbf{F}_2)_{\std}$ is generated as a $\mathbf{F}_2[x_1,x_2,x_3]^{S_3}$-module by homogeneous copies of $\std$, so following~\cite{wang2022explicit}, we consider \textit{generating representations} of $Q_m(3,\mathbf{F}_2)_\std$ as homogeneous copies of $\std$ in a generators and relations presentation of $Q_m(3,\mathbf{F}_2)_{\std}$ with a minimal generator set.

\subsubsection{Quasi-invariants at half-integers}

Note that if $\k$ is a field with $\mathrm{char\,}\k\neq 2$ and $m\in\Z_{\geq 0}$, then for any $K\in\k[x_1,\dots,x_n]$, $(x_{i_1}-x_{i_2})^{2m}|(1-s_{i_1i_2})K$ implies $(x_{i_1}-x_{i_2})^{2m+1}|(1-s_{i_1i_2})K$ since $(1-s_{i_1i_2})K$ is $s_{i_1i_2}$-antiinvariant, hence the exponent $2m+1$ in the definition of quasi-invariant polynomials. But this does not hold in characteristic 2, since there is no concept of antiinvariants. Indeed, one can check that for $K=x_1^2+x_2^2$, we have $(x_{i_1}-x_{i_2})^2|(1-s_{i_1i_2})K$ for all $i_1$, $i_2$, but $(x_{i_1}-x_{i_2})^3\nmid|(1-s_{i_1i_2})K$ if~$i_1=1,2$,~$i_2\neq 1,2$.

We encapsulate this data by extending the definition of quasi-invariants to half-integers when ${p=2}$. For example, $K=x_1^2+x_2^2$ is \smash{$\frac{1}{2}$}-quasi-invariant, and this is in fact the minimal degree nonsymmetric \smash{$\frac{1}{2}$}-quasi-invariant polynomial. Proposition~\ref{firstProp} still holds when $m,m'$ are half-integers, and the definitions of $Q_m(3,\mathbf{F}_2)_\triv$, $Q_m(3,\mathbf{F}_2)_\std$ also naturally extend to half-integer~$m$. So from now on, whenever we refer to quasi-invariants in characteristic $2$ we let $m$ be a~half-integer.

